\begin{itemize}
    \item RNF1.1: Pertinencia funcional (Adecuación funcional). Es la capacidad del producto software para proporcionar un conjunto de funciones que facilitan la consecución de tareas y objetivos de usuario especificados.

    \item RNF2.1: Auto-descriptividad (Capacidad de interacción). Es la capacidad de un producto para presentar la información adecuada, haciendo su uso inmediatamente evidente para el usuario sin interacciones excesivas con el producto u otros recursos (documentación, \textit{help desks}, otros usuarios).

    \item RNF3.1: Capacidad para ser probado (Mantenibilidad). Es la facilidad con la que se pueden establecer criterios de prueba para un sistema o componente y con la que se pueden llevar a cabo las pruebas para determinar si se cumplen dichos criterios.

    \item RNF3.2: Modularidad (Mantenibilidad). Es la capacidad de un producto para evitar que los cambios en un componente afecten a otros componentes.
\end{itemize}
